\ficha{Weller (2018), Sovereign Debt Crises and Negotiations}{weller2018}

\begin{identificacao}
WELLER, Leonardo. \textit{Sovereign Debt Crises and Negotiations in Brazil and
Mexico, 1888-1914: Governments versus Bankers}. Cham: Palgrave Macmillan, 2018.
178 p. DOI 10.1007/978-3-319-73633-4.

Livro, monografia derivada de tese de doutorado na London School of Economics,
em inglês. Lidos por inteiro os capítulos 1 (p. 3-20), 2 (p. 21-39), 4
(Rothschilds' Troubled Republic: Brazil, 1889-1898, p. 59-81) e 5 (Rothschilds
and Coffee Finance: Brazil, 1898-1914, p. 83-102), mais a conclusão (p. 163-167).
Os capítulos 3, 6, 7 e 8 foram apenas consultados.
\end{identificacao}

\bloco{Problema e tese}

Quem manda na negociação de um empréstimo soberano, o banco emissor ou o governo
devedor? Weller responde que o governo era jogador independente, e que a
literatura sobre a primeira globalização financeira erra ao tratar o lado da
demanda como passivo. Governos escolhiam entre conceder monopólio a um banco
patrono e negociar com bancos concorrentes, e essa escolha, não a vontade do
banqueiro, explica as condições de crédito obtidas.

\bloco{Argumento}

O livro inteiro sustenta uma tese comparativa. Brasil e México fizeram escolhas
opostas: o Brasil entregou aos Rothschild o monopólio de sua dívida de 1850 a
1908, enquanto o México nunca emitiu dois empréstimos consecutivos com o mesmo
grupo de bancos (p. 6-7). Cada opção tinha preço. O banco patrono emprestava a
marca, abria conexões de negócios e oferecia guarda-chuva de crédito na crise; a
concorrência baixava comissão e elevava o preço de emissão (p. 22, 30). A
comparação com o México acrescenta o custo do caminho oposto: sem banco patrono
exposto, o governo revolucionário não obteve empréstimo de socorro e a dívida
entrou em moratória em 1914 (p. 163-166).

O quadro analítico do capítulo 2 tem dois eixos. Governo é forte quando tem bom
histórico de pagamento, estabilidade política e contas sólidas; banco é forte
quando tem status alto e pouca exposição ao cliente. O achado contraintuitivo
está no cruzamento: quando um cliente exclusivo bate no teto de crédito, o status
premier torna o banco fraco, não forte, porque a moratória destruiria sua marca
(p. 36). Daí o conceito central, o \textit{rescue loan}, crédito emitido a prêmio
de risco abaixo da taxa de mercado durante a crise (p. 11, 27). O socorro é o
seguro que o cliente paga ao aceitar a exclusividade. O padrão-ouro entra apenas
como literatura resenhada: Weller resume Bordo, para quem a adesão era ``good
housekeeping seal of approval'' (p. 24), e contrapõe Flandreau e Zumer, para quem
os fundamentos fiscais, e não o regime cambial, determinavam o custo do serviço
da dívida (p. 25).

O capítulo 4 refaz a crise dos anos 1890 como caso de carona. Rui Barbosa,
influenciado pela teoria quantitativa do século XIX, expandiu a emissão e o M2
cresceu 144,9\% entre 1890 e 1891 (p. 63). Os Rothschild subscreveram 13,1
milhões de libras entre 1893 e 1895, e o empréstimo de 1893 saiu a 90\% do risco
Brasil, o que o qualifica como socorro (p. 67). O contrafactual mostra que, sem o
aumento do gasto militar ou sem a transferência do imposto de exportação aos
estados, o governo teria pago o serviço da dívida com recursos próprios (p. 70).
O funding loan de 1898 não foi imposição: Prudente de Morais recusou a oferta
condicionada ao arrendamento da Central e mandou insinuar a suspensão de
pagamentos para forçar acordo maior (p. 73-74). O contrato saiu com 8,6 milhões
de libras, três anos de juros pagos em títulos e suspensão de amortização até
1911, com a queima de papel fiscalizada pelo London and River Plate Bank (p. 75).
A taxa foi negociada: os banqueiros queriam queimar papel até 16 pence,
Bernardino de Campos obteve 18, e o M2 caiu 31\% em vez de 34\% (p. 76). Weller
chama isso de propriedade compartilhada da política e atribui aos Rothschild o
papel de corretor de poder que fortaleceu Murtinho contra o lobby do café
(p. 76-78).

O capítulo 5 é o mais próximo do objeto da dissertação. A resposta do café à
valorização do mil-réis foi dupla: estoque e câmbio. São Paulo tomou 1 milhão de
libras do Disconto Gesellschaft em 1906 e mais 3 milhões via Henry Schroeder,
enquanto Rodrigues Alves recusava aval federal e os Rothschild chamavam o esquema
de ``impossible speculation'' e ``most dangerous and suicidal policy'' (p. 89).
Com os estoques em 8 milhões de sacas no fim de 1907, o governo federal socorreu
São Paulo por fora, com cerca de 1,2 milhão de libras retiradas da conta que
mantinha com os próprios Rothschild, cujo saldo caiu de 6,1 milhões em 1906 para
1,7 milhão em 1908 (p. 90). O banco não quis se identificar com o esquema, mas
participou do socorro, porque uma moratória dos bônus do café derrubaria a
cotação de todos os títulos brasileiros (p. 91).

O segundo braço é a Caixa de Conversão, que Weller trata como peça da política do
café. O governo depositava o câmbio que entrava e emitia moeda contra ele a 15
pence por mil-réis, abaixo da taxa de mercado, que orbitava 16 em 1907 (p. 91).
Daí o adjetivo do título da seção: o padrão-ouro brasileiro foi montado para
depreciar, não para valorizar a moeda, e, ao contrário do que a literatura sobre
o padrão clássico faria esperar, permitiu ao governo imprimir moeda, com o M2
subindo já em 1907 (p. 91-92). A consequência é uma política pró-cíclica, com o
saldo externo favorável financiando défice e absorção de dívida interna (p. 93).
Em 1910 a Caixa acumulou 20 milhões de libras e a paridade foi valorizada de 15
para 16 pence (p. 93-94), enquanto o défice do governo Hermes da Fonseca saltava
de 15\% da receita tributária em 1909 para 81\% em 1914 (p. 93). Com a queda do
café, o M2 encolheu 13\% entre o início de 1913 e meados de 1914, quando o
governo abandonou a Caixa e deixou o câmbio depreciar (p. 94).

A quebra do monopólio aparece como cálculo de custo, não como conflito. O governo
emitiu em Paris por Société Générale, Paribas, Crédit Mobilier e Caisse
Commerciale para pagar comissão menor: 0,75\% na emissão e 0,5\% no resgate no
contrato de 1911, contra 4,25\% cobrados pelos Rothschild em 1908 (p. 86-87). Os
banqueiros reclamaram em privado e não retaliaram (p. 87-88). Em 1914 ainda
socorreram o cliente infiel: recusada pelo governo a proposta de abril, que
exigia penhor das rendas alfandegárias e arrendamento da Central (p. 99), o
sindicato ofereceu depois da eclosão da guerra até 15 milhões de libras a 1,67\%
de prêmio de risco, cerca de metade do risco Brasil, com treze anos sem
amortização e sem condição alguma de política econômica, porque os banqueiros
entenderam a crise como cíclica e não estrutural (p. 96, 99-100).

\chave{Brazil's Gold Standard was peculiar: It was sat up to depreciate (rather
than to appreciate) the national currency}{p. 91}

\chave{in the 1910s an unusual Gold Standard had allowed the government to run
lose economic policy until coffee prices collapsed}{p. 96}

\bloco{Conceitos e definições operacionais}

\textit{Rescue loan}: crédito concedido durante crise a prêmio de risco inferior
ao risco do país no mercado secundário (p. 11, 27); distingue-se de
\textit{settlement}, posterior à moratória, e de \textit{rescue package}, que soma
socorro e recontratação (p. 27-28). \textit{Risco do país}: diferença entre o
rendimento dos títulos do país em Londres e o do consol britânico (p. 14).
\textit{Exposição}: reputacional, não financeira, medida pela participação do
cliente na carteira de bônus com a marca do banco (p. 68). Padrão-ouro
brasileiro: paridade fixa administrada por uma \textit{caisse} que funciona como
currency board, mas calibrada abaixo do mercado e compatível com expansão
monetária (p. 91-92).

\bloco{Evidência e método}

Correspondência do Rothschild Archive (séries XI e I), Biblioteca Histórica do
Ministério da Fazenda (contratos e o dossiê 332/532 do funding loan), Instituto
Histórico e Geográfico Brasileiro (cartas de Prudente de Morais a Bernardino de
Campos), Arquivo Nacional (Coleção Afonso Pena) e arquivos mexicanos. Séries de
prêmio de risco construídas a partir de \textit{The Investor's Monthly Manual};
dados fiscais e monetários do IBGE (1990) e do Balanço da receita e despesa da
República; imprensa financeira britânica como termômetro de expectativa. Método:
narrativa de arquivo disciplinada pelo quadro de dois eixos, com contrafactuais
aritméticos (p. 70) e correlações entre risco brasileiro e índice de emergentes
(p. 86, 94).

\bloco{Agentes e vertentes}

Na teoria: Eaton e Gersovitz, Tomz, Mitchener e Weidenmier (sanções), Bordo
(padrão-ouro como selo), Flandreau e Zumer (fundamentos fiscais), Flandreau e
Flores (marcas e bancos como porteiros), Esteves e Tunçer. No Brasil: Rui Barbosa
como expansionista de fundamentação quantitativista (p. 63); Rodrigues Alves,
Prudente de Morais e Bernardino de Campos na negociação de 1898; Campos Sales e
Joaquim Murtinho como núcleo ortodoxo (p. 77); David Campista como destinatário
das cartas dos banqueiros sobre o café (p. 89); Afonso Pena autorizando o socorro
a São Paulo (p. 90); Hermes da Fonseca como marco da deterioração fiscal (p. 93);
Francisco Sales e Rivadávia Corrêa em 1913 e 1914. Nos bancos: Alfred e Gustave
Rothschild, Disconto Gesellschaft, Schroeder, Paribas, Société Générale.

\bloco{Críticas e limites}

Weller declara que deixa de fora a interação entre bancos e portadores finais
(p. 28), o que enfraquece a afirmação de que os Rothschild convenciam o mercado.
O livro é história de negociação, não de política monetária interna: a Caixa de
Conversão aparece como instrumento do lobby do café, sem discussão do debate
legislativo, das taxas alternativas ou do limite de emissão. O empréstimo de
valorização de 1908 em Paris não é tratado; os empréstimos franceses que o livro
registra entre 1908 e 1911 são federais e destinados a ferrovias e conversão
(p. 12, 86). A caracterização do padrão-ouro brasileiro como heterodoxo apoia-se
em Peláez e Suzigan e em Fritsch, sem confronto com a literatura que lê a Caixa
como currency board ortodoxo.

\begin{dialogo}
\item Procurar a defesa explícita da taxa baixa como proteção ao café, e não como
etapa rumo ao par legal. Se o argumento aparecer nesses termos em O Paiz e no
Correio Paulistano, confirma a tese de que a Caixa foi montada para depreciar
(p. 91) e enfraquece a leitura da adesão ao padrão-ouro como busca de
credibilidade externa.
\item Rastrear como a imprensa liga a Caixa ao esquema de valorização em 1907 e
1908. Weller documenta socorro federal a São Paulo com recursos da conta dos
Rothschild (p. 90) sem menção pública; ausência ou presença do tema nos jornais
testa se a operação era segredo de gabinete.
\item Verificar se as emissões em Paris foram noticiadas como afirmação de
soberania financeira ou como assunto técnico de comissão. Weller as lê como
cálculo de custo (p. 86-87), não como ruptura com os Rothschild.
\item Procurar, em 1913 e 1914, o vocabulário de moratória e de novo funding.
Weller mostra o governo usando a ameaça de suspender pagamentos como alavanca
(p. 97-99). Se os jornais discutem a ameaça abertamente, o debate público era
parte da negociação, não seu reflexo.
\item Conferir se aparece a crítica de que a Caixa permitia expansão monetária e
défice, núcleo do diagnóstico de Weller sobre 1908-1913 (p. 93).
\end{dialogo}

\usonocapitulo{Parte II, como peça central. Serve à afirmação de que o
enforcement das instituições ortodoxas sobre o Brasil não foi imposição
unilateral do banqueiro, e sim resultado de barganha em que a exposição
reputacional dos Rothschild dava alavancagem ao governo. Sustenta também a ponte
para a Parte IV: a Caixa de Conversão de 1906 é, na leitura de Weller, um
padrão-ouro calibrado para depreciar, montado sob pressão do lobby do café e
compatível com emissão, o que dá base à hipótese de que o debate de 1906 gira em
torno do nível da taxa e da proteção ao café, e não da conversibilidade em si.}
